% !TeX root = ../main.tex

\chapter*{Abstract}
\addcontentsline{toc}{chapter}{Abstract}
Tijdens deze industriële stage van heb ik voor MotorSport Training Center een lokale
desktopapplicatie ontwikkeld die live timingdata omzet in bruikbare informatie voor tijdens de race. De bestaande timingwebsites tonen de officiële rangschikking en actuele tijden, maar
bieden niet alle teamgerichte vergelijkingen, voorspellingen en strategische hulpmiddelen die
tijdens een endurancewedstrijd nodig zijn. De ontwikkelde Electron-applicatie leest daarom periodiek
de timingpagina, normaliseert data van GetRaceResults en RIS Timing, bewaart de racehistoriek lokaal
en berekent analyses buiten de gebruikersinterface.

Het resultaat is een dashboard voor maximaal drie gelijktijdig gevolgde wagens, met afzonderlijke
modi voor race, kwalificatie en vrije training. Het systeem berekent onder andere ronde- en
sectorgemiddelden, geldige beste tijden, voorspelde rondetijden, marges tegenover configureerbare
referentietijden, vergelijkingen tussen piloten en klassewagens, inhaalindicaties en een
pitstopplanning. Neutralisaties, pitgerelateerde rondes en onvolledige data worden expliciet
behandeld. Een afzonderlijk grafiekvenster maakt langere trends zichtbaar. De data wordt per sessie
opgeslagen in leesbare JSON-, JSONL- en CSV-bestanden, waardoor een sessie na een onderbreking kan
worden hervat.

De ontwikkeling verliep iteratief en in nauw overleg met de opdrachtgever. Praktijktests en
screenshots brachten verschillende subtiele fouten aan het licht, onder meer bij rondeachterstanden,
intervalvelden, pitstopprojecties en wisselende timingproviders. Deze fouten leidden tot modulaire
domeinfuncties en een uitgebreide automatische testsuite. Builds voor Windows en macOS worden via
GitHub Actions gemaakt en als release gepubliceerd. Een praktijktest tijdens een raceweekend in Spa
bevestigde dat de volledige dataketen met RIS Timing operationeel werkte, maar bracht bijkomende
tekortkomingen aan het licht bij vlagovergangen, selectie van geldige pace-ronden en stabiliteit van
gaps. De opgeslagen Spa-data vormt de basis voor de volgende regressietests en verbeteringen.

\textbf{Trefwoorden:} motorsport, live timing, Electron, race-engineering, rondetijdanalyse,
pitstopstrategie, softwaretesten.
